\RequirePackage{bmpsize-base}
\RequirePackage{ifpdf}
\makeatletter
% 中译本改动：这段 DVI 模式下手工计算图片包围盒的 hack，是为真正的
% "dvips -> ps2pdf" 流程写的，假设的是 pdftex 的原语。ifpdf 包在 XeLaTeX/
% LuaLaTeX 下都可能判定成"不是 pdf 模式"从而误触发这段 hack，报
% "Undefined control sequence"（真实复现：\@bmpsize@read 报错，编译直接
% 中止；用 \ifdefined\XeTeXversion/\luatexversion 分别探测两种引擎又在
% LuaLaTeX 下报 "\ifdefined ... was incomplete"，怀疑是这两个宏在当前
% texlive 版本下的可用性问题，没有深究）。中译本只会用 XeLaTeX 或
% LuaLaTeX 编译，两者都能自己正确处理图片包围盒，不需要这段 hack，直接
% 无条件跳过最省事也最不容易出岔子。
\expandafter\@gobble{%
\AtBeginDocument{%
\let\Gin@ii@old=\Gin@ii
\def\Gin@ii[#1]#2{%
  \begingroup
    \let\@found\@empty
    \@for\@type:=\bmpsize@types\do{%
      \ifx\@found\@empty
        \@nameuse{bmpsize@read@\@type}{#2.\@type}%
        \ifbmpsize@ok
          \let\@found=\@type
        \fi
      \fi
      \ifx\@found\@empty
        \@nameuse{bmpsize@read@\@type}{#2}%
        \ifbmpsize@ok
          \let\@found=\@type
        \fi
      \fi
    }%
    \ifx\@found\@empty
      \Gin@ii@old[#1]{#2}%
    \else
      \Gin@ii@old[natwidth=\bmpsize@width bp,natheight=\bmpsize@height bp,#1]{#2}%
    \fi
  \endgroup
}%
}
}
\makeatother
